\chapter{Pensamiento Computacional y Algoritmos}

\section{¿Qué es un algoritmo?}

Un \textbf{algoritmo} es una secuencia finita de pasos ordenados y precisos que describe cómo resolver un problema. Para que una secuencia de pasos sea considerada un algoritmo debe cumplir las siguientes características:

\begin{itemize}
    \item \textbf{Finito:} debe terminar en algún momento, no puede ejecutarse indefinidamente.
    \item \textbf{Preciso:} cada paso debe ser claro e inequívoco, sin ambigüedades.
    \item \textbf{Definido:} si se ejecuta dos veces con los mismos datos de entrada, siempre produce el mismo resultado.
    \item \textbf{Con entrada y salida:} puede recibir datos del exterior (entradas) y debe producir al menos un resultado (salida).
\end{itemize}

\begin{notabox}
Los algoritmos son independientes del lenguaje de programación. Se pueden expresar en lenguaje natural, en pseudocódigo o en un diagrama de flujo, y luego traducirse a cualquier lenguaje de programación.
\end{notabox}

\textbf{Ejemplo cotidiano --- receta de cocina:}
\begin{enumerate}
    \item Calentar agua hasta que hierva.
    \item Agregar los fideos.
    \item Esperar 8 minutos.
    \item Colar y servir.
\end{enumerate}

Este es un algoritmo: tiene pasos ordenados, es finito y produce un resultado determinado.

\section{Representación de algoritmos: el diagrama de flujo}

Un \textbf{diagrama de flujo} es una representación gráfica de un algoritmo. Utiliza símbolos estandarizados conectados por flechas que indican el orden de ejecución de los pasos.

\subsection{Símbolos básicos}

\tikzCircuito{cap4/simbolos-flujo}{Símbolos básicos del diagrama de flujo: forma de cada uno y para qué se usa.}

\subsection{Reglas generales}

\begin{itemize}
    \item Todo diagrama empieza con un símbolo \textbf{Inicio} y termina con uno o más símbolos \textbf{Fin}.
    \item Las flechas indican la dirección del flujo: normalmente de arriba hacia abajo y de izquierda a derecha.
    \item Un rombo de decisión siempre tiene exactamente \textbf{dos salidas}: una para cuando la condición es verdadera (Sí) y otra para cuando es falsa (No).
    \item No deben existir caminos que queden sin salida ni flechas que apunten ``al vacío''.
\end{itemize}

\section{Estructuras de control}

Todo algoritmo se construye combinando tres estructuras fundamentales:

\subsection{Secuencia}

Los pasos se ejecutan uno tras otro, en el orden en que aparecen. Es la estructura más simple.

\tikzCircuito{cap4/secuencia}{Estructura de secuencia: cálculo del área de un rectángulo.}

\textbf{Ejemplo:} calcular el área de un rectángulo.
\begin{enumerate}
    \item Leer base.
    \item Leer altura.
    \item Calcular área = base $\times$ altura.
    \item Mostrar área.
\end{enumerate}

\subsection{Decisión (selección)}

El flujo se divide en dos caminos según el resultado de una condición. Se evalúa una expresión lógica (verdadera o falsa) y se ejecuta uno u otro bloque de pasos.

\tikzCircuito{cap4/si-sino}{Estructura de decisión (si--sino): las ramas Sí y No se reúnen al final.}

\textbf{Ejemplo:} determinar si un número es positivo o negativo.

\begin{lstlisting}[style=estilo-pseudocode]
Si número > 0 entonces
    Mostrar "El número es positivo"
Sino
    Mostrar "El número es negativo o cero"
Fin Si
\end{lstlisting}

Las decisiones pueden encadenarse o anidarse para evaluar múltiples condiciones.

\subsection{Repetición (iteración)}

Un bloque de pasos se ejecuta más de una vez, mientras se cumpla una condición. Hay dos variantes principales:

\subsubsection{Repetir mientras (while --- condición al inicio)}

La condición se evalúa \textbf{antes} de ejecutar el bloque. Si la condición es falsa desde el principio, el bloque no se ejecuta nunca.

\tikzCircuito{cap4/mientras}{Estructura de repetición \emph{mientras}: la condición se evalúa antes del bloque.}

\begin{lstlisting}[style=estilo-pseudocode]
Mientras condición sea verdadera hacer
    [bloque de pasos]
Fin Mientras
\end{lstlisting}

\textbf{Ejemplo:} mostrar los números del 1 al 10.
\begin{lstlisting}[style=estilo-pseudocode]
contador = 1
Mientras contador <= 10 hacer
    Mostrar contador
    contador = contador + 1
Fin Mientras
\end{lstlisting}

\subsubsection{Repetir hasta (do-while --- condición al final)}

La condición se evalúa \textbf{después} de ejecutar el bloque, por lo que el bloque se ejecuta \textbf{al menos una vez}. En PSeInt, el bucle se repite \emph{mientras la condición sea falsa} y termina \emph{cuando la condición se vuelve verdadera}: se repite ``hasta que'' se cumpla.

\tikzCircuito{cap4/repetir-hasta}{Estructura de repetición \emph{repetir--hasta que}: primero se ejecuta el bloque y después se evalúa la condición.}

\begin{lstlisting}[style=estilo-pseudocode]
Repetir
    [bloque de pasos]
Hasta Que condición
\end{lstlisting}

\textbf{Ejemplo:} pedir un número hasta que sea positivo.
\begin{lstlisting}[style=estilo-pseudocode]
Repetir
    Leer numero
Hasta Que numero > 0
\end{lstlisting}

\begin{notabox}
\textbf{Mientras vs.\ repetir--hasta:} comparando los dos diagramas se ve la diferencia. En el \emph{mientras} el rombo está \textbf{antes} del bloque: si la condición es falsa de entrada, el bloque no se ejecuta nunca. En el \emph{repetir--hasta} el rombo está \textbf{después}: el bloque siempre se ejecuta al menos una vez. Además, el \emph{mientras} sigue repitiendo con \textbf{Sí}, y el \emph{repetir--hasta} sigue repitiendo con \textbf{No}.
\end{notabox}

\section{Variables}

Una \textbf{variable} es un espacio de memoria con nombre que almacena un valor que puede cambiar durante la ejecución del algoritmo.

Características:
\begin{itemize}
    \item Tiene un \textbf{nombre} (identificador) que la distingue de las demás.
    \item Almacena un \textbf{tipo de dato} (número entero, número decimal, texto, etc.).
    \item Su valor puede \textbf{leerse} (consultar qué tiene) o \textbf{escribirse} (asignarle un nuevo valor).
\end{itemize}

\subsection{Tipos de datos básicos}

\begin{table}[H]
\centering
\begin{tabular}{lp{5cm}l}
\toprule
Tipo & Descripción & Ejemplo \\
\midrule
Entero & Número sin decimales & 5, -3, 100 \\
Real & Número con decimales & 3.14, -0.5 \\
Carácter & Un símbolo de texto & 'A', '7' \\
Cadena & Texto (varios caracteres) & "Hola", "temperatura" \\
Lógico & Verdadero o falso & Verdadero, Falso \\
\bottomrule
\end{tabular}
\end{table}

\subsection{Asignación}

La asignación guarda un valor en una variable. En pseudocódigo se escribe con la flecha $\leftarrow$:

\begin{lstlisting}[style=estilo-pseudocode]
temperatura <- 25
resultado <- base * altura
nombre <- "Juan"
\end{lstlisting}

\section{Operadores}

\subsection{Aritméticos}

\begin{table}[H]
\centering
\begin{tabular}{cl}
\toprule
Operador & Operación \\
\midrule
\texttt{+} & Suma \\
\texttt{-} & Resta \\
\texttt{*} & Multiplicación \\
\texttt{/} & División \\
\texttt{MOD} & Módulo (resto de la división entera) \\
\bottomrule
\end{tabular}
\end{table}

\subsection{Relacionales (comparación)}

\begin{table}[H]
\centering
\begin{tabular}{cl}
\toprule
Operador & Significado \\
\midrule
\texttt{>} & Mayor que \\
\texttt{<} & Menor que \\
\texttt{>=} & Mayor o igual que \\
\texttt{<=} & Menor o igual que \\
\texttt{=} & Igual \\
\texttt{<>} & Distinto \\
\bottomrule
\end{tabular}
\end{table}

\subsection{Lógicos}

\begin{table}[H]
\centering
\begin{tabular}{cp{9cm}}
\toprule
Operador & Significado \\
\midrule
\texttt{Y} / \texttt{AND} & Verdadero solo si ambas condiciones son verdaderas \\
\texttt{O} / \texttt{OR} & Verdadero si al menos una condición es verdadera \\
\texttt{NO} / \texttt{NOT} & Invierte el valor lógico \\
\bottomrule
\end{tabular}
\end{table}

\section{Herramienta: PSeInt}

\textbf{PSeInt} es un programa que permite escribir y ejecutar algoritmos en pseudocódigo, y visualizarlos como diagrama de flujo de forma automática.

\subsection{Configuración del perfil}

Para este curso se usa el perfil \textbf{EEST1 VL}, disponible en el aula virtual (Moodle). Para cargarlo:

\begin{enumerate}
    \item Ir a \textbf{Configurar $\rightarrow$ Opciones del lenguaje (perfiles) $\rightarrow$ Cargar}.
    \item Seleccionar el archivo \texttt{EEST1\_VL}.
    \item Hacer clic en \textbf{Aceptar}.
\end{enumerate}

\subsection{Editar el diagrama de flujo}

Existen tres formas de acceder al modo de diagrama de flujo:

\begin{itemize}
    \item Hacer clic en el ícono de diagrama de flujo en la barra de menú.
    \item Ir a \textbf{Archivo $\rightarrow$ Editar diagrama de flujo}.
    \item Presionar \textbf{F7}.
\end{itemize}

\subsection{Ejecutar el algoritmo}

\begin{itemize}
    \item \textbf{Ejecución completa:} botón \textit{Ejecutar} (flecha verde), o \textbf{F9}.
    \item \textbf{Ejecución paso a paso:} botón \textit{Paso a paso} (ícono de pies), útil para seguir el flujo y detectar errores.
\end{itemize}

\section{Ejercicios}

Los siguientes ejercicios deben resolverse primero dibujando el diagrama de flujo en PSeInt y verificando su funcionamiento con distintos valores de entrada.

\subsection{Ejercicio 1 --- Suma de dos valores}

Realizar un algoritmo capaz de sumar 2 valores e informar el resultado.

\subsection{Ejercicio 2 --- Operación con tres valores}

Realizar un algoritmo que permita ingresar 3 valores, sumar los primeros 2, restarle el tercero al resultado, e informar el resultado final.

\textbf{Relación esperada:} resultado = (A + B) $-$ C

\subsection{Ejercicio 3 --- Módulo de un valor}

Realizar un algoritmo capaz de calcular el módulo (valor absoluto) de un número e informar el resultado. Recordar que el módulo devuelve el valor positivo: si el número es positivo lo devuelve igual, si es negativo lo invierte.

\begin{advertenciabox}
Este ejercicio introduce la estructura de decisión: se necesita evaluar si el número es negativo para invertirlo.
\end{advertenciabox}

\subsection{Ejercicio 4 --- Sensor de temperatura}

Realizar un algoritmo capaz de informar si hace calor o frío. Se ingresa un valor de temperatura y el algoritmo debe:

\begin{itemize}
    \item Si la temperatura supera los 26°C: informar ``hace calor''.
    \item Si la temperatura es menor a 22°C: informar ``hace frío''.
    \item En cualquier otro caso (entre 22° y 26°): no informar nada (temperatura agradable).
\end{itemize}

\begin{advertenciabox}
Este ejercicio introduce decisiones múltiples: más de una condición debe evaluarse.
\end{advertenciabox}

\subsection{Ejercicio 5 --- Promedio de una serie de valores}

Realizar un algoritmo capaz de calcular el promedio de todos los números ingresados. El ingreso termina cuando el usuario introduce un 0.

\begin{advertenciabox}
Este ejercicio introduce la estructura de repetición. El algoritmo no sabe de antemano cuántos números se van a ingresar: debe seguir pidiendo valores hasta que se ingrese un 0, y recién entonces calcular y mostrar el promedio.
\end{advertenciabox}

\textbf{Consideraciones:}
\begin{itemize}
    \item Llevar la cuenta de cuántos valores se ingresaron (contador).
    \item Llevar la suma acumulada de los valores.
    \item No incluir el 0 en el cálculo del promedio.
    \item Tener en cuenta el caso en que se ingrese 0 como primer valor (no hay valores para promediar).
\end{itemize}
